% crossterm.tex -- round 476, PRIME.RDAGGER.CROSSTERM_SHARPENING.01.
% UNCONDITIONAL on the first Dirichlet H^2 ball through L_16=2.7726.
% NOT lambda_*(L)>=0.  scramble-sensitive.
% Companion: verify_crossterm.py.
% Discovery: experiments/tfpt-discovery/crossterm_probe.py.
% Fixed L only.  NO RH CLAIM.  NO anti-RH claim.
\documentclass[11pt]{article}

\usepackage[a4paper,margin=2.45cm]{geometry}
\usepackage{amsmath,amssymb,amsthm}
\usepackage{booktabs,array,microtype,xcolor}
\usepackage[T1]{fontenc}
\usepackage[colorlinks=true,linkcolor=blue!50!black,urlcolor=blue!50!black]{hyperref}

\newtheorem{lemma}{Lemma}
\newtheorem{theorem}{Theorem}
\theoremstyle{remark}
\newtheorem{remark}{Remark}
\newcommand{\QW}{Q_{\mathrm W}}
\setlength{\emergencystretch}{2em}

\title{\bfseries Archimedean cross-term sharpening\\[2mm]
\large PRIME.RDAGGER.CROSSTERM\_SHARPENING.01 --- round 476}
\author{Stefan Hamann}
\date{August 31, 2026}

\begin{document}
\maketitle

\begin{abstract}\noindent
The r474 diagnosis is implemented at fixed $L$: the slack was the
archimedean cross-term, not the grid.  The first Dirichlet $H^2$
ball with $\sigma_2=(\pi/(2L))^2$ is one-dimensional, spanned by
$h_0(x)=\cos(\pi x/(2L))$.  Route~B evaluates $\QW(h_0)/\|h_0\|_2^2$
from the closed autocorrelation, the r475 regularized $u$-space
pairing, a Bose prefix, and the exact algebraic tail
$-L\bigl(\psi(K+1)-\psi(K+\tfrac14)\bigr)$.  The Yoshida/Bombieri
gate $L=0.3$ (prime term identically zero) returns
$r\in[1.163\cdot10^{-2},\,1.180\cdot10^{-2}]>0$.  Every sealed
schedule point through $L_{16}=2.7726$ is positive on this ball.
This is not $\lambda_*(L)\ge0$.  Verdict
\texttt{UNCONDITIONAL($L_{\max}=2.7726$)} on the declared class.
No RH claim.
\end{abstract}

\medskip\noindent
\textbf{Status.}\enspace
\texttt{UNCONDITIONAL($L_{\max}=2.7726$)} on the first Dirichlet
$H^2$ ball.  No promotion of an RH claim.

\begin{center}
\fbox{\parbox{0.95\linewidth}{%
\textbf{Claim boundary.}  Nothing here is $\lambda_*(L)\ge0$ on the
r461 $L^2$ class.  The theorem is the Rayleigh number of a single
smooth generator.  The archimedean symbol $\sigma_A(0)<0$, so $A$
is not a positive operator; $\Pi$ supplies the reserve at low
frequency.  No $L^*$ claim, no $R^\dagger$ claim, no RH claim,
no anti-RH claim.}}
\end{center}

\section{Declared class}

After r474, the working comparison class is
\[
 \mathcal B(L)=\bigl\{h\in H^2_0([-L,L]):
 \|h''\|_2\le(\pi/(2L))^2\|h\|_2\bigr\}.
\]
The first Dirichlet eigenvalue of $-d^2/dx^2$ on $[-L,L]$ with
vanishing endpoints is $(\pi/(2L))^2$.  Hence
$\mathcal B(L)=\mathrm{span}\{h_0\}$,
$h_0(x)=\cos(\pi x/(2L))$.  No silent $L^2$ enlargement.

\section{Closed data}

Write $\omega=\pi/(2L)$.  The autocorrelation $g=h_0*\widetilde h_0$
is even, supported in $[-2L,2L]$, and
\[
 g(t)=\bigl(L-\tfrac t2\bigr)\cos(\omega t)+\frac{\sin(\omega t)}{2\omega}
 \qquad(t\in[0,2L]),
\]
with $g(0)=\|h_0\|_2^2=L$.  The r471 form is
$\QW=A-P+\Pi$, scramble-sensitive:
$P(g)=\sum_{n\le e^{2L}}2\Lambda(n)n^{-1/2}g(\log n)$.
At $L=0.3$ one has $2L<\log2$, so $P\equiv0$.

\begin{lemma}[Regularized arch pairing]\label{lem:A}
With $w(u)=2e^{-u/2}/(1-e^{-2u})$ and
$C_b=-(\gamma+\log\pi)-\log(1-e^{-4L})$,
\[
 A(g)=C_b\,g(0)+\int_0^{2L}w(u)\bigl(e^{-3u/2}g(0)-g(u)\bigr)\,du.
\]
The integrand is regular at $u=0$.  Expanding
$w(u)=2\sum_{k\ge0}e^{-(2k+1/2)u}$ and integrating termwise against
the closed $g$ yields a Bose prefix of $K=200$ terms plus the exact
algebraic tail $-L\bigl(\psi(K+1)-\psi(K+\tfrac14)\bigr)$ and an
$O(K^{-2})$ cubic remainder.
\end{lemma}

\begin{lemma}[Pole, closed]\label{lem:Pi}
$\Pi(g)=2\bigl(I(-\tfrac12)+I(\tfrac12)\bigr)$ with
$I(\alpha)=\int_0^{2L}e^{-\alpha u}g(u)\,du$ elementary.
\end{lemma}

\begin{lemma}[Symbol at zero]\label{lem:sig}
The Fourier symbol of $A$ obeys
$\sigma_A(0)=-(\gamma+\log\pi)-(\pi/2+3\log2)\in[-5.37219,-5.37218]$.
In particular $A$ is negative at zero frequency.  The series
$\sigma_A(t)=-(\gamma+\log\pi)+\sum_k 2\bigl[(2k+2)^{-1}
-\alpha_k/(\alpha_k^2+t^2)\bigr]$
meets the closed value after the same $\psi$-tail.
\end{lemma}

\begin{lemma}[Peano constant of the regularized kernel]\label{lem:Peano}
$\displaystyle C_0:=\int_0^\infty w(s)\frac{s^2}{2}\,ds
=\sum_{k\ge0}2\alpha_k^{-3}$
with $\alpha_k=2k+\tfrac12$, and
$C_0\in[16.16596,\,16.16597]$.
Against an even increment with
$|\Delta g(s)-\Delta g(0)|\le(s^2/2)\|\Delta g''\|_\infty$
this is the $L^1$ factor of the regularized pairing (PROVED
constant; not used for the Rayleigh table).
\end{lemma}

\section{Theorem}

\begin{theorem}[First-mode unconditional]\label{thm:main}
On $\mathcal B(L)$ one has
$\QW(h)/\|h\|_2^2=r(L)$ with $r(L)$ the sealed interval of
Section~\ref{sec:table}.  In particular $r(L)>0$ at every
schedule point, including the Yoshida/Bombieri value $L=0.3$
and $L_{16}=4\log2=2.7726$.
\end{theorem}

The combination implication of r474 is not needed: there is no
interpolant and no $\omega_L(\delta)$ on this one-dimensional ball.

\section{Sealed balance}\label{sec:table}

Probe \texttt{crossterm\_probe.py},
SPEC\_SHA \texttt{a123650e757b912d}.
Smoke $15/15$ in $0.14\,\mathrm{s}$; full $25/25$ in $0.79\,\mathrm{s}$.
Every row is \texttt{H2BALL\_POS}.  YB-gate PASS.

\begin{center}
\small
\begin{tabular}{@{}lrrrl@{}}
\toprule
$L$ & $r_{\mathrm{lo}}$ & $r_{\mathrm{hi}}$ & $\#\{p^k\}$ & verdict\\
\midrule
$0.30$ (YB) & $1.163\cdot10^{-2}$ & $1.180\cdot10^{-2}$ & $0$ & H2BALL\_POS\\
$0.80$ & $1.837\cdot10^{-4}$ & $2.080\cdot10^{-4}$ & $3$ & H2BALL\_POS\\
$L_5=0.866$ & $9.209\cdot10^{-4}$ & $9.416\cdot10^{-4}$ & $4$ & H2BALL\_POS\\
$1.0$ & $1.099\cdot10^{-4}$ & $1.254\cdot10^{-4}$ & $5$ & H2BALL\_POS\\
$1.2$ & $8.729\cdot10^{-5}$ & $9.808\cdot10^{-5}$ & $8$ & H2BALL\_POS\\
$L_8=1.386$ & $1.389\cdot10^{-4}$ & $1.470\cdot10^{-4}$ & $10$ & H2BALL\_POS\\
$1.4$ & $8.724\cdot10^{-5}$ & $9.517\cdot10^{-5}$ & $10$ & H2BALL\_POS\\
$1.6$ & $9.875\cdot10^{-5}$ & $1.048\cdot10^{-4}$ & $13$ & H2BALL\_POS\\
$1.8$ & $9.679\cdot10^{-5}$ & $1.016\cdot10^{-4}$ & $18$ & H2BALL\_POS\\
$2.0$ & $7.191\cdot10^{-5}$ & $7.579\cdot10^{-5}$ & $24$ & H2BALL\_POS\\
$L_{12}=2.079$ & $2.456\cdot10^{-5}$ & $2.816\cdot10^{-5}$ & $27$ & H2BALL\_POS\\
$L_{16}=2.773$ & $3.452\cdot10^{-6}$ & $5.473\cdot10^{-6}$ & $70$ & H2BALL\_POS\\
\bottomrule
\end{tabular}
\end{center}

At $L=0.3$, $A\approx-0.28954$ and $\Pi\approx+0.29305$.  The pole
reserve beats the negative well of $\sigma_A$.  The $L_{16}$
margin is thin but strictly positive.

\section{What this does not do}

It does not prove $\lambda_*(L)\ge0$ at any $L$.  It does not
restate a cofinal quantifier.  High modes and functions that fail
$h(\pm L)=0$ remain outside $\mathcal B(L)$.  No RH claim.

\end{document}
